% ---------------------------------------------------------------------------
% OPTION A, CHAPTER 2. Written 2026-11-10 with AI assistance (Claude) from
% docs/thesis/OUTLINE_BRIEF.md (chapter 2, candidate sources) and
% CHAPTER_MATERIAL.md on branch claude/thesis-brief-verification-6573e8.
% Eleven sources were added to references.bib on 2026-11-10 with metadata
% from Crossref, arXiv and CEUR-WS. Of the two preprints in the comparison
% table, Slide2Text (arXiv:2503.17710) was read in full; for Learn Your Way
% (arXiv:2509.13348) the abstract and the study design were read, not every
% section. The table compares what the papers describe; nothing was run.
% Figure F2 (positioning) is no longer used: it was drawn from this system's
% own capabilities, not from a literature search.
% The author must read, revise and own every sentence before submission.
% ---------------------------------------------------------------------------
\chapter{Background and Related Work}
\label{ch:background}

This chapter gives the background that the three main chapters rely on and
places Learnstation among related systems. It follows the path of a page:
how software recovers text from a PDF, what a language model adds and where it fails, how a system ties answers to a source, what the learning sciences say about practice, and how a system can return data about learners without exposing them.
The chapter ends with a comparison of the two systems closest to ours.

\section{Recovering Content from Documents}
\label{sec:document-ai}

A PDF describes where glyphs and images sit on a page. It does not
say which glyphs form a heading, which a table, or in which order a human
reads them. Document understanding is the field that recovers this structure.
Models such as LayoutLM learn from text and position
together~\cite{xu2020layoutlm}, data sets such as DocLayNet provide
human-annotated page layouts to train and test them
on~\cite{pfitzmann2022doclaynet}, and Nougat transcribes scientific pages
directly from the image into markup~\cite{blecher2023nougat}. Open-source
toolkits such as Docling~\cite{auer2024docling} and
MinerU~\cite{wang2024mineru} package layout analysis, table recognition, and
optical character recognition into one converter, and benchmarks such as
OmniDocBench compare converters across document
types~\cite{ouyang2025omnidocbench}. A simpler path reads the text layer that
a born-digital PDF already carries, for which we use
PyMuPDF~\cite{pymupdf}.

Lecture slides are a hard case for all of these, for a reason that is not
technical. A slide holds little text, and much of its meaning is in a figure,
an arrow, or the position of two boxes. SlideVQA shows that questions about slide decks can need more than one slide and its visual content~\cite{tanaka2023slidevqa}. Dinh et al.\ compare ways to give a
language model access to university course material and find that passing
retrieved slides as images works better than passing their extracted
text~\cite{dinh2025slides}. \Cref{ch:content} reaches a related result by another route: in a blind comparison, the author preferred the text that a vision model wrote from the page image over every text extractor.
None of the works above compares extractors on lecture slides by what a
learner would rather read, or at the cost a small server can bear.
\Cref{ch:content} does both.

\section{Language Models and Their Failure Modes}
\label{sec:llms}

A large language model produces fluent text for any prompt, whether or not
the text is true. The literature calls a fluent and unsupported statement a
hallucination and treats it as a property of the technique and not a rare
fault~\cite{ji2023hallucination,huang2025hallucination}. For learning
material this matters more than elsewhere, because a learner is by definition
not in a position to notice the error.

Two further failure modes shape \cref{ch:trust}. The first concerns the use
of a model to judge the output of a model. Such judges agree with human
raters to a useful degree, but they carry biases of
their own~\cite{zheng2023judging}, and a model tends to rate its own output
more favourably than that of others~\cite{panickssery2024selfpreference}.
This is why the answer check in \cref{ch:trust} runs on a different model
than the one that wrote the question. The second is prompt injection. When a
model reads a document, instructions hidden in the document compete with the
instructions of the application~\cite{greshake2023indirect,perez2022ignore},
and OWASP lists the risk first among the risks of LLM applications~\cite{owasp2025llm}. Marking untrusted text as data, called
spotlighting, lowers the success of such attacks without removing
them~\cite{hines2024spotlighting}. Every uploaded PDF is untrusted text in
this sense.

\section{Retrieval and Groundedness}
\label{sec:rag}

Retrieval-augmented generation answers a question in two steps. It first
retrieves passages from a corpus and then lets the model write an answer from
them~\cite{lewis2020rag,gao2023ragsurvey}. Retrieval commonly compares dense vectors of question and passage~\cite{karpukhin2020dpr}, in some systems combined with keyword search, and an approximate index such as HNSW (hierarchical navigable small-world
graphs) keeps the comparison
fast~\cite{malkov2020hnsw}. We store the vectors in PostgreSQL with the
pgvector extension~\cite{pgvector}.

Retrieval makes an answer checkable, and it does not make it correct. Work on
attribution separates two questions: whether a citation points at a source,
and whether the source supports the sentence that cites
it~\cite{rashkin2023ais}. Models that a prompt asks to cite often fail the second~\cite{gao2023alce}, and metrics such as FActScore~\cite{min2023factscore} and RAGAS~\cite{es2024ragas} measure that gap. We use the word \emph{grounded} for an answer that a system restricts to retrieved material and the word \emph{verified} only where a
separate step has checked the content.

A system can enforce groundedness at different strengths.
\Cref{fig:groundedness} shows the four that exist in Learnstation, from a gate that refuses before any model call to a prompt that asks the model
to stay with the page. \Cref{ch:trust} describes each of them and how the
system moved between them.

\thesisfigure{f3-grounded-vs-ungrounded}
  {Four groundedness regimes over the same material. The course tutor behind
   the ask bar of a Space refuses before any model call. The system checks Luna's page references and highlighted phrases against the page. Explain
   and Ask in the reading view rely on the prompt alone, and the tutor of the
   older interface, which has no refusal gate, no longer reaches students. \Cref{ch:trust} introduces each of them.}
  {fig:groundedness}

\section{What the Learning Sciences Ask For}
\label{sec:learning-science}

Results from the learning sciences explain why the system does more
than display a page. They motivate the practice features of \cref{ch:system}, which this thesis does not evaluate, and they explain why the assistant stays close to the page (\cref{ch:trust}). First, retrieving something from memory strengthens it
more than reading it again~\cite{roediger2006testing}, and practice spread
over time beats the same practice in one block~\cite{cepeda2006distributed}.
A review of ten study techniques rates these two as the most
useful~\cite{dunlosky2013techniques}. Schedulers in the SM-2 family turn the
spacing result into an algorithm that sets the next review of an item from
the learner's last answers~\cite{wozniak1990sm2}. The scheduler behind Luna's
practice cards is a variant of it.

Second, doing goes with more learning than reading or watching.
Koedinger et al.\ found this in the data of an online course and are careful
to call it a correlation~\cite{koedinger2015spectator}. Van Campenhout et
al.\ place automatically generated questions into textbooks to obtain such
practice at scale~\cite{vancampenhout2021transforming}, and reviews of
automatic question generation describe how hard it is to judge the quality
of the result~\cite{kurdi2020aqg}. Third, the same content in more than one
form can support understanding when the forms complement each other. Our reading view, which shows a page as the original and as texts written
from it, is
consistent with this line of
work~\cite{ainsworth1999functions,ainsworth2006deft}, but we do not claim that we derived it from this work.

The evidence on AI tutors is mixed, and the mix is instructive. Human
tutoring and intelligent tutoring systems reach similar
effects~\cite{vanlehn2011tutoring}. A tutor built on a language model beat
active learning in class in one controlled
study~\cite{kestin2025aitutoring}. In another, students who practised with an
unrestricted chat model did worse once they lost access to it, and a version with guardrails largely avoided that harm~\cite{bastani2025guardrails}. A conversational tutor that a model generated from textbook text improved how learners
perceived the session and did not improve what they learned compared to
reading~\cite{schmucker2024ruffle}. We read these results as support for a
design in which the assistant stays with the material and makes the learner
do something, and as a warning against claiming a learning effect without
measuring it. Game elements such as points and badges, which the system also
has, show small positive effects on average and depend on the
context~\cite{deterding2011gamification,hamari2014gamification,sailer2020gamification}.

\section{Feeding Data Back Without Exposing Learners}
\label{sec:analytics-privacy}

Learning analytics collects traces of what learners do and returns them to
someone who can act on them, and Clow describes this as a cycle that is only
complete when the action reaches the learners again~\cite{clow2012cycle}.
The person who can act on a confusing page is the lecturer who made it. The
traces are personal data, however, and in a course of twenty the lecturer
knows every reader. The General Data Protection Regulation asks for data protection by
design and by default~\cite{gdpr2016}. A classic technical measure is a
minimum group size: under $k$-anonymity a released record must be
indistinguishable from at least $k-1$ others~\cite{sweeney2002kanonymity}.
A minimum group size has known limits. It does not protect a value that
every member of the group shares, which $\ell$-diversity
addresses~\cite{machanavajjhala2007ldiversity}, and it says nothing about
what two releases reveal together. Differential privacy limits what any number of releases reveals about one
person, at a privacy cost that adds up over releases~\cite{dwork2006dp}.
\Cref{ch:loop} applies the idea of a minimum group size not to released records but to released statements: the view says nothing about a page that fewer than five people have
read. We call a view privacy-preserving in this thesis if no statement it
releases is about fewer than five people. \Cref{ch:loop} measures the
owner's view against this criterion.

\section{Related Systems}
\label{sec:positioning}

Intelligent textbooks are the closest tradition. Brusilovsky et al.\ describe
their return: textbooks that model their content, embed practice, and record
use, now with machine learning for the steps that once required
authors~\cite{brusilovsky2022return}. A second line augments the reading
surface of a document without changing the document. Paper Plain adds plain
language summaries and guiding questions to medical papers, and its authors evaluated it with 24~participants against a standard PDF reader~\cite{august2023paperplain}. The Semantic Reader project builds such augmented reading interfaces for
scientific papers in general~\cite{lo2024semanticreader}. Learnstation
shares the principle of this line: the original stays visible, and the
generated layer sits on top of it.

Two recent systems transform teaching material with language models and are
close enough to compare directly (\cref{tab:related}). Learn Your Way, from
a team at Google, turns a textbook into several personalised representations
``while maintaining content integrity'' and reports a randomised controlled
trial with 60~students aged 15 to 18 against a standard PDF
reader~\cite{learnlm2025textbook}. The treatment in that trial includes
quizzes that the control does not have, so the trial does not isolate the
effect of the reading surface. Slide2Text takes PowerPoint files and writes a
textbook from them~\cite{zhou2025slide2text}. It enriches the slides with
passages retrieved from external papers and web search, which is the opposite
of our choice to stay inside the uploaded material. Both are preprints. We
did not run either system, and the table compares what the papers describe.

\begin{table}[htbp]
  \centering\small
  \caption{Learnstation and the two closest systems, compared on what their
    publications describe. We ran no system for this comparison.}
  \label{tab:related}
  \begin{tabular}{@{}p{0.16\textwidth}p{0.245\textwidth}p{0.245\textwidth}p{0.245\textwidth}@{}}
    \toprule
     & \textbf{Learn Your Way}~\cite{learnlm2025textbook}
     & \textbf{Slide2Text}~\cite{zhou2025slide2text}
     & \textbf{Learnstation} \\
    \midrule
    Input &
      Textbook material &
      PowerPoint files (text, notes, images by OCR) &
      Any lecture PDF: slides, scripts, exercise sheets \\
    Output &
      Several representations of the same content &
      A textbook with chapters, cases, and exercises &
      The original page, a text of it, questions, and a search index \\
    Relation to the source &
      Generated from the source, content integrity as a stated goal &
      Source enriched with retrieved external papers and web results &
      Original page stays the default view; answers cite pages; the ask bar
      refuses outside the material \\
    Checks on generated content &
      Ratings by pedagogical experts &
      A review by several models and an expert panel, described as a
      framework &
      Computed answer keys, and an independent check on a different model \\
    Personalisation &
      Yes, of the content itself &
      Language, style, and difficulty chosen by the user &
      None of the content; only the schedule of practice \\
    Feedback to the owner &
      None described &
      None described &
      Patterns per page from five readers on, no names \\
    Evaluation &
      Controlled trial with 60~students &
      One case study on a deck of 39~slides; a controlled study is planned,
      no results reported &
      Blind benchmark, scripted readers, a probe of 16~pairs; no learners \\
    \bottomrule
  \end{tabular}
\end{table}

The comparison shows where this thesis stands. Both related systems aim at
the individual learner and personalise the content. As far as
the publications describe, neither has a path back to the owner of the material and neither rejects a
generated item automatically. Learnstation makes the opposite trade: every
reader gets the same content, the owner gets a view of how people read it, and a check runs on generated quiz questions before a reader sees them, unless the
check itself is unavailable. What Learn Your
Way has and this thesis lacks is evidence from real learners, and
\cref{ch:conclusion} names its trial as the model for the study that should
follow. The comparison with commercial tools in the January proposal, which
contrasted the project with quiz and presentation products, was the two developers' own reading of public product pages, and we do not repeat it here.
